\subsection{Section summary}\label{subsec:abc-section-summary}

This section described how agreement-based cascading, the cost-aware routing system, works and how it relates to the MAP stage of the knowledge-extraction pipeline. ABC replaces a single MAP call with a three-tier voter ensemble, and at each tier, the voter outputs are compared using a max-consensus agreement scorer. Embedding-only and hybrid strategies can be used as the similarity strategy of the agreement scorer. Based on the calculated deferral score, and the calibrated acceptance and rejection thresholds, the cascade decides whether to accept the voter outputs at a certain tier, reject them, or escalate to the next tier with more capable but expensive voters. 

In this section, we also explained the safety policies of the cascade: when there is a conflict in polarity direction between voters in a tier on the same subject-outcome-relation-category combination, the cascade is forced to escalate. There is also a single-voter policy, which defines how the cascade should behave in case there is only one eligible voter output within a tier. These policies are intended to avoid structurally unsafe accept decisions by ABC; their purpose is not to improve strict F1 scores.

The thresholds, models and flag configurations in the cascade are calibrated by an offline replay mechanism, where all voter outputs are cached, removing API calls from the calibration loop, which speeds up the calibration procedure and reduces costs. 

Whether the agreement-based cascade actually improves the cost-quality trade-off on the histopathology corpus is an empirical question, answered in Chapter~\ref{chapter:Results} and interpreted in Chapter~\ref{chapter:Discussion}.